\documentclass[literature_review.tex]{subfiles}

\begin{document}
The converse of the monolithic kernel design is a microkernel which operates on Liedtke's principle of minimality \textit{`A concept is tolerated inside the $\mu$-kernel only if moving it outside the kernel, i.e., permitting competing implementations, would prevent the implementation of the system's required functionality'} \citep{Liedtke_95}. Essentially if a component is not strictly required to be inside the kernel for the system to function, then it should be implemented outside the kernel, i.e. in userland. \TODO{add a diagram here for microkernel.}

\subsection{Microkernel Security}
The argument for microkernels revolves around security. The idea of a verified \gls{tcb} or kernel is ideal as that is the starting point to building a truly trustworthy system. The microkernel design allows for the verification of the entire kernel as exemplified by seL4, the world's first verified microkernel \citep{Klein_AEHCDEEKNSTW_10}. Whereas it is theoretically possible to formally verify a monolithic operating system, the inherently massive size of the \gls{tcb} of a monolithic kernel already makes the cost and effort of doing so prohibitive. Proof maintenance is also required as the codebase changes according to changing requirements. The fact that only the strictly necessary components of a microkernel are included in the microkernel means that changes are less frequent and less significant. Conversely, a monolithic operating system will have frequent and often large scale changes to its \gls{tcb} already evidenced by Linux's \gls{tcb}'s fast rate of growth over the last three decades meaning that proof maintenance costs will be extremely high. Another factor is that while the difficulty of formally reverifying a component if it is local and low-level is relatively low and proportional to the size of the change, the cost of adding new, large, cross-cutting features is exponentially more expensive \citep{Klein_AEHCDEEKNSTW_10}. Due to the highly coupled nature of monolithic operating systems, the verification of monolithic and even hybrid operating systems can be assumed to be unfeasible. 

A study into how a multiserver operating system based on a verified microkernel can eliminate security vulnerabilities by analysing all critical security bugs from the \textit{Common Vulnerabilities and Exposures} theorises that 96\% of all critical Linux compromises would be not critical with a microkernel based design, 40\% would be completely eliminated on a verified microkernel and 29\% on an unverified microkernel \citep{Biggs_LH_18}. This shows just how much of an improvement kernel design choices can have and that even without verification, \gls{mkbs}'s are just inherently more secure. The reason for this is that principled implemenation of a microkernel based system inherently induces a design based on the principle of \textit{least privilege} \citep{Saltzer_Schroeder_75} as services can be given only the bare minimum authority to do their job. Therefore when said service is compromised, only the privileges held by the service already are affected. By the same principle, microkernels promote a design which enforces fault containment as only the components that depend on a compromised service will be affected \citep{Biggs_LH_18}.

\subsection{Performance of Microkernels}
As compared to monolithic systems, \gls{mkbs}'s take on an inherent cost when servicing a system call. During a syscall on a \gls{mkbs} at least four mode switches and two context switches is required: Client invokes kernel's IPC operating (2 mode switches) and so does the server (2 mode switches). The kernel in delivering the messages switches context from client to server and back. Meanwhile, the same syscall on a monolithic system only requires two mode switches (one to trap into kernel mode and another to resume execution in usermode) \citep{Biggs_LH_18}. Hence, with the implementation of a performant \gls{mkbs} it is a question of how much of this cost is amortised by other procedures and how much can we optimise the kernel and the operating system servers of the \gls{mkbs} given that the microkernel codebase is magnitudes smaller and the individual components more modular and easier to reason about.

\subsection{Early Microkernel Concept} 
The concept of the microkernel was first defined in Brinch Hansen's work on the RC 4000 multiprogramming system \textit{The Nucleus of a Multiprogramming System} which outlines a system substrate including proccess management including synchronous \gls{ipc}, process creation, control \citep{Brinch_Hansen_70}. This \textit{nucleus} allowed one to build their own system around it, implementing one's own policies as required for the task, much like the role a modern microkernel would play in a microkernel based operating system. This generality is what is innovative about \textit{nucleus} as it does not make any assumptions about the workload of the user, instead allowing the user to extend the \textit{nucleus} to create a system tailored to their requirements through the minimal interface which the \textit{nucleus provides}.

Hansen outlines the concept of processes as an execution of one or more programs in a storage area, with the ability to start, stop and remove processes and also communicate synchronously with other processes, with the \textit{nucleus} facilitating these actions. Processes would form a tree, the parent being able to control (start, stop or remove) its children and its children's descendants. In such system, the \textit{nucleus} would be analogous to a microkernel and an \textit{operating system} would be the root of the hierarchy tree.

\subsection{First Generation Microkernels}
Mach 3.0 is an example of one of the earliest microkernels.
\TODO{go into how the concept of a microkernel has developed since nucleus.}
\TODO{talk about mach microkernel}
\TODO{talk about other microkernels}

\subsection{Second Generation Microkernels}
Learning from the mistakes of the first generation microkernels.
\TODO{talk about L34}
\subsection{Third Generation Microkernels}
\TODO{talk about L4Re}
\TODO{talk about other 3gen microkernels}
\subsection{General Purpose Microkernels}
\subsubsection{Mach 3.0}
\TODO{talk about mach3.0}

\subsubsection{HongMeng}
An example of a state of the art microkernel based general purpose operating system is HongMeng. HongMeng is unique in the fact that it is a truly commercialised general purpose microkernel that seems to avoid the problems which seemed to plague IBM's workplace \gls{os} and Chorus. However, it is important to note that HongMeng's desing introduces many choices which are a departure from Liedtke's minimality principle.

\paragraph{Differentiated Isolation Classes} are a concept introduced in HongMeng where instead of placing all \gls{os} services in userland, \gls{os} services are differntiated by performance criticality and trust and placed into \textit{isolation classes}. There are 3 isolation classes: IC0, IC1 and IC2. IC0 applies to the core \gls{os} services and is described as the \textit{core \gls{tcb}}. In deployment, the Linux ABI-compliant shim is included in IC0. \gls{ipc} calls between IC0 are indirect function calls. IC1 is still within the kernel space and isolation between components is enforced by mechanisms as different services are assigned domains where ARM watchpoint and Intel PKS prevents cross domain memory accesses. \gls{ipc} between IC1 or to IC0 enters a gate in the microkernel which performs a minimal context switch. It is important to note that IC1 and IC0 can be considered the real \gls{tcb} as this is the code that runs in a privileged mode. IC2 is what is the userland in a traditional microkernel and \gls{ipc} incurs the same cost as a normal \gls{ipc}.

In extension to placing certain \gls{os} services essentially back inside the kernel, HongMeng also explores the idea of coalesced \gls{os} services in order to reduce \gls{ipc}. This idea is analogous to coarse grained servers for microkernel services that sacrifice security for performance.

\paragraph{Address Token-based Access Control} is an alternative security measure to capabilities which is what traditional \gls{sota} microkernels like seL4 use. The claim is that accessing the kernel objects behind capabilities are slow requiring privilege switches and accesses to multiple metadata tables. By mapping kernel objects to an \gls{os} service's vspace, we can grant permissions at a lower cost, especially when multiple kernel objects are batched into one mapping. However, the disadvantage is that the restricting of fine grained operations and chain revocations are impossible due to implicit ownership. The compromise is that HongMeng retains usage of capabilities and uses address tokens only for certain kernel objects. The potential problem with adding such functionality is that the size of the kernel increases, greatly increasing the kernel's cache footprint - one of the main reasons for Mach 3.0's poor performance \citep{Chen_Bershad_93, Liedtke_96b}.

\paragraph{Policy Free Kernel Paging} is implemented in HongMeng contrary to a userland pager as an additional round trip from the kernel to the pager induces a significant performance degradation. Hence, policy driven decisions are made in advance allowing an in kernel page fault handling mechanism. This results in a page fault latency that is comparable to Linux and much faster than L4Re, while only compromising on the ability to change the policy after being preallocated to the PCache with some negative impacts to the PCache.

\paragraph{Linux ABI Shim} is a key inclusion in the implementation of HongMeng. One of the main challenges which microkernels face when attempting to `go general' is the fact that Linux has carried the standard for \gls{abi}'s, hence software is created for Linux's \gls{abi}. Since the reimplementation of all software for a microkernel's native \gls{abi}, certain workarounds must be made. The Linux \gls{abi} shim is a component in IC0 that redirects Linux syscalls into \gls{ipc}'s towards their respective \gls{os} servers. Since it resides in IC0, the additional overhead of an extra \gls{ipc} is avoided. This component also acts as a central repository of state such as file descriptors, allowing for an easier implementation of fork and poll which are syscalls that are thought of as challenging to implement.

\paragraph{Driver Containers} are used in order to reuse Linux drivers without modification by exposing all necessary k\gls{api}'s to the driver. The `container' holds the Linux drivers which can be invoked through \gls{vfs} and the `container' is assigned physical devices. However as the Linux Driver Container is placed in IC2, in order to reduce latency for performance critical drivers, a `twin driver' is created in the native driver container which is in IC1 to handle \gls{io} \gls{irq}'s.  

\subsubsection{Djawula}
Djawula is a prototype of a proof of concept general purpose \gls{os} based on seL4. The aim of this project is to explore the feasibility of a verifiable general purpose operating system that is designed with adherence to the principle of minimality in mind. It can be seen as a logical step forward after LionsOS \citep{Heiser_VCJGNLDZAZWPDB_25} to see whether similar performance can be maintained even with a larger project scope. The main difference between LionsOS and Djawula is that while resource allocation and \gls{pd}'s are statically configured on compile time in LionsOS, Djawula would have dynamic resource allocation and the ability to load \gls{pd}'s in runtime which significantly increases complexity. The development of Djawula is in early stages.

Djawula currently consists as a collection of cores servers: boot file server, init server and root server. The boot file server is the first process and this server spawns the init server and then enters a system call handler loop, redirecting syscalls to the rootserver or to \gls{sddf} components \citep{Heiser_CBDP_24}. The init server spawns the \gls{sddf} components and the rootserver is currently a coarse grained server that handles all of the core \gls{os} services. In future iterations, the rootserver should be divided into a collection of fine grained servers to improve cache effects and allow for easier modification, policy changes and verification.
\end{document}